
\documentclass{article}
\usepackage{colortbl}
\usepackage{makecell}
\usepackage{multirow}
\usepackage{supertabular}

\begin{document}

\newcounter{utterance}

\centering \large Interaction Transcript for game `hot\_air\_balloon', experiment `air\_balloon\_survival\_en\_negotiation\_hard', episode 5 with Qwen3{-}32B{-}t0.0.
\vspace{24pt}

{ \footnotesize  \setcounter{utterance}{1}
\setlength{\tabcolsep}{0pt}
\begin{supertabular}{c@{$\;$}|p{.15\linewidth}@{}p{.15\linewidth}p{.15\linewidth}p{.15\linewidth}p{.15\linewidth}p{.15\linewidth}}
    \# & $\;$A & \multicolumn{4}{c}{Game Master} & $\;\:$B\\
    \hline

    \theutterance \stepcounter{utterance}  
    & & \multicolumn{4}{p{0.6\linewidth}}{
        \cellcolor[rgb]{0.9,0.9,0.9}{
            \makecell[{{p{\linewidth}}}]{
                \texttt{\tiny{[P1$\langle$GM]}}
                \texttt{You are participating in a collaborative negotiation game.} \\
\\ 
\texttt{Together with another participant, you must agree on a single set of items that will be kept. Each of you has your own view of how much each item matters to you (importance). You do not know how the other participant values the items. Additionally, you are given the effort each item demands.} \\
\texttt{You may only agree on a set if the total effort of the selected items does not exceed a shared limit:} \\
\\ 
\texttt{LIMIT = 3386} \\
\\ 
\texttt{Here are the individual item effort values:} \\
\\ 
\texttt{Item effort = \{"C17": 313, "B41": 720, "A71": 175, "C00": 461, "B35": 635, "A77": 685, "A03": 544, "A04": 203, "C54": 369, "C01": 539, "A87": 4, "B32": 695, "B86": 399, "B33": 594, "A10": 437\}} \\
\\ 
\texttt{Here is your personal view on the importance of each item:} \\
\\ 
\texttt{Item importance values = \{"C17": 138, "B41": 583, "A71": 65, "C00": 262, "B35": 121, "A77": 508, "A03": 461, "A04": 484, "C54": 668, "C01": 389, "A87": 215, "B32": 97, "B86": 500, "B33": 30, "A10": 400\}} \\
\\ 
\texttt{Goal:} \\
\\ 
\texttt{Your goal is to negotiate a shared set of items that benefits you as much as possible (i.e., maximizes total importance to YOU), while staying within the LIMIT. You are not required to make a PROPOSAL in every message {-} you can simply negotiate as well. All tactics are allowed!} \\
\\ 
\texttt{Interaction Protocol:} \\
\\ 
\texttt{You may only use the following structured formats in a message:} \\
\\ 
\texttt{PROPOSAL: \{'A', 'B', 'C', …\}} \\
\texttt{Propose keeping exactly those items.} \\
\texttt{REFUSE: \{'A', 'B', 'C', …\}} \\
\texttt{Explicitly reject opponent's proposal.} \\
\texttt{ARGUMENT: \{'...'\}} \\
\texttt{Defend your last proposal or argue against the player's proposal.} \\
\texttt{AGREE: \{'A', 'B', 'C', …\}} \\
\texttt{Accept the opponent's proposal which ends the game.} \\
\texttt{STRATEGIC REASONING: \{'...'\}} \\
\texttt{	Describe your strategic reasoning or anticipation explaining your choice of action. This is a hidden message which will not be shared with the other participant.} \\
\\ 
\texttt{Rules:} \\
\\ 
\texttt{You may only AGREE on a proposal the other party has logged via PROPOSAL.} \\
\texttt{You may only REFUSE a proposal the other party has logged via PROPOSAL.} \\
\texttt{Total effort of any PROPOSAL or AGREE set must be ≤ LIMIT.} \\
\texttt{Do NOT reveal your hidden importance scores.} \\
\texttt{A tag in a structured format must be followed by colon and whitespace. The argument must be a python set containing 0 or more strings.} \\
\texttt{So, it must be of the form TAG: \{...\}} \\
\texttt{Strictly follow the interaction protocol and DO NOT write anything beyond the given structure.} \\
\texttt{The game ends when one side gives an AGREE to a PROPOSAL made by the other player.} \\
\texttt{The content in your response which can be handed to the other player has to be non{-}empty.} \\
\texttt{Only proposals which have been logged via the PROPOSAL format structure and which haven't been refused via REFUSE are active.} \\
\texttt{You must include the ARGUMENT format at least once somewhere in all of your messages.} \\
\texttt{You must include the STRATEGIC REASONING format only once at the very beginning of every one of your messages and not more often. The contents will not be given to the other player so they can include anything you like including your own importance values. Here you should reason step by step to come up with you next move.} \\
            }
        }
    }
    & \\ \\

    \theutterance \stepcounter{utterance}  
    & \multicolumn{4}{p{0.6\linewidth}}{
        \cellcolor[rgb]{0.9,0.9,0.9}{
            \makecell[{{p{\linewidth}}}]{
                \texttt{\tiny{[P1$\rangle$GM]}}
                \texttt{<think>} \\
\texttt{Okay, let's see. I need to figure out the best set of items to propose that maximizes my total importance without exceeding the effort limit of 3386. First, I should calculate the total effort for different combinations. But since I don't know the other person's preferences, I have to prioritize my own high{-}importance items while keeping the effort low.} \\
\\ 
\texttt{Looking at my importance values, the highest ones are C54 (668), B41 (583), A77 (508), B86 (500), C17 (138), A03 (461), A04 (484), A87 (215), A10 (400), etc. Wait, maybe I should sort them by importance per unit effort to see which items give me the most value for the effort they take.} \\
\\ 
\texttt{Let me calculate the importance per effort ratio for each item:} \\
\\ 
\texttt{{-} C54: 668 / 369 ≈ 1.81} \\
\texttt{{-} B41: 583 / 720 ≈ 0.81} \\
\texttt{{-} A77: 508 / 685 ≈ 0.74} \\
\texttt{{-} B86: 500 / 399 ≈ 1.25} \\
\texttt{{-} A04: 484 / 203 ≈ 2.38} \\
\texttt{{-} A03: 461 / 544 ≈ 0.85} \\
\texttt{{-} A10: 400 / 437 ≈ 0.92} \\
\texttt{{-} A87: 215 / 4 = 53.75 (Wow, that's super high!)} \\
\texttt{{-} C00: 262 / 461 ≈ 0.57} \\
\texttt{{-} C01: 389 / 539 ≈ 0.72} \\
\texttt{{-} B35: 121 / 635 ≈ 0.19} \\
\texttt{{-} B32: 97 / 695 ≈ 0.14} \\
\texttt{{-} B33: 30 / 594 ≈ 0.05} \\
\texttt{{-} A71: 65 /} \\
            }
        }
    }
    & & \\ \\

    \theutterance \stepcounter{utterance}  
    & & \multicolumn{4}{p{0.6\linewidth}}{
        \cellcolor[rgb]{0.9,0.9,0.9}{
            \makecell[{{p{\linewidth}}}]{
                \texttt{\tiny{[P1$\langle$GM]}}
                \texttt{Your response did not start with the proper strategic reasoning tag at the very beginning of your response. The very first structured format must be of the form STRATEGIC REASONING: \{...\}. Try again.} \\
            }
        }
    }
    & \\ \\

    \theutterance \stepcounter{utterance}  
    & \multicolumn{4}{p{0.6\linewidth}}{
        \cellcolor[rgb]{0.9,0.9,0.9}{
            \makecell[{{p{\linewidth}}}]{
                \texttt{\tiny{[P1$\rangle$GM]}}
                \texttt{<think>} \\
\texttt{Okay, let me start over. The user wants me to make sure the first part is the STRATEGIC REASONING. I need to outline my thought process here.} \\
\\ 
\texttt{First, I need to maximize my total importance without exceeding the effort limit. Let me list out the items sorted by importance per effort ratio. The highest ratios are A87 (215/4=53.75), A04 (484/203≈2.38), C54 (668/369≈1.81), B86 (500/399≈1.25), then others. So I should prioritize these. } \\
\\ 
\texttt{A87 has the highest ratio and very low effort. Including it is a must. Then A04, which is also high ratio. C54 and B86 next. Let's check the total effort if I take A87 (4), A04 (203), C54 (369), B86 (399). Sum so far: 4+203=207, +369=576, +399=975. That's under 3386. Now, adding other high importance items. B41 (720 effort, 583 importance) would add 720, total 975+720=1695. Still under. Then A77 (685 effort, 508 importance) adds 685 to 1695+685=2380. Next, A03 (544, 461) would take total to 2380+544=2924. Then A10 (437, 400) adds 437 to 2924+437=3361. Total effort is 3361, which is under 3386. That leaves 25 effort left. } \\
\\ 
\texttt{Now, check if there are other items that can fit. C00 (461) would exceed if added. C17 (313) would add 313 to 3361=3674, which is over} \\
            }
        }
    }
    & & \\ \\

    \theutterance \stepcounter{utterance}  
    & & \multicolumn{4}{p{0.6\linewidth}}{
        \cellcolor[rgb]{0.9,0.9,0.9}{
            \makecell[{{p{\linewidth}}}]{
                \texttt{\tiny{[P1$\langle$GM]}}
                \texttt{Your response did not start with the proper strategic reasoning tag at the very beginning of your response. The very first structured format must be of the form STRATEGIC REASONING: \{...\}. Try again.} \\
            }
        }
    }
    & \\ \\

    \theutterance \stepcounter{utterance}  
    & \multicolumn{4}{p{0.6\linewidth}}{
        \cellcolor[rgb]{0.9,0.9,0.9}{
            \makecell[{{p{\linewidth}}}]{
                \texttt{\tiny{[P1$\rangle$GM]}}
                \texttt{<think>} \\
\texttt{Okay, I need to make sure I start with the STRATEGIC REASONING tag. Let me outline my plan again.} \\
\\ 
\texttt{First, I'll calculate the importance per effort for each item to prioritize those that give me the most value for the least effort. The highest ratios are A87 (215/4=53.75), A04 (484/203≈2.38), C54 (668/369≈1.81), B86 (500/399≈1.25), and others. I should include A87 first because it's super efficient. Then add A04, C54, B86, and so on, checking the total effort each time to stay under 3386.} \\
\\ 
\texttt{Next, I'll list the items in order of priority and sum their efforts. Starting with A87 (4), A04 (203), C54 (369), B86 (399), B41 (720), A77 (685), A03 (544), A10 (437). Let's add these up: 4+203=207, +369=576, +399=975, +720=1695, +685=2380, +544=2924, +437=3361. Total effort is 3361, which is under the limit. That leaves 25 effort units. Maybe add C17 (313) would exceed, so skip. The next item is C00 (461) which is too much. So the current set is optimal for my importance.} \\
\\ 
\texttt{Now, I need to propose this set. The total importance would be A87(215) + A04(484) + C54(668) + B86(500) + B41(583) + A77(508) + A03(461) + A10(400) = Let's calculate: 215+484=699, +6} \\
            }
        }
    }
    & & \\ \\

    \theutterance \stepcounter{utterance}  
    & & & \multicolumn{2}{p{0.3\linewidth}}{
        \cellcolor[rgb]{0.9,0.9,0.9}{
            \makecell[{{p{\linewidth}}}]{
                \texttt{\tiny{[GM$|$GM]}}
                \texttt{abort} \\
            }
        }
    }
    & & \\ \\

    \theutterance \stepcounter{utterance}  
    & & & \multicolumn{2}{p{0.3\linewidth}}{
        \cellcolor[rgb]{0.9,0.9,0.9}{
            \makecell[{{p{\linewidth}}}]{
                \texttt{\tiny{[GM$|$GM]}}
                \texttt{end game} \\
            }
        }
    }
    & & \\ \\

\end{supertabular}
}

\end{document}
